\documentclass[11pt]{article}
\usepackage[margin=1in]{geometry}
\usepackage{amsmath,amsthm,amssymb,mathtools}
\usepackage[colorlinks=true,linkcolor=blue,citecolor=blue,urlcolor=blue]{hyperref}

\newtheorem{proposition}{Proposition}
\newtheorem{theorem}[proposition]{Theorem}
\newtheorem{corollary}[proposition]{Corollary}
\newtheorem{remark}[proposition]{Remark}

\newcommand{\torus}[2]{D_{#1}(#2)}
\newcommand{\cycle}[1]{\vec C_{#1}}

\title{A Product Reduction for Composite Dimensions}
\author{}
\date{}

\begin{document}
\maketitle

\begin{abstract}
This note isolates a theorem-level reduction for composite dimensions in the
directed torus Hamilton-decomposition problem. The key point is that the class
of dimensions admitting Hamilton decompositions for every modulus is
multiplicatively closed. Consequently, once the prime dimensions are settled,
all composite dimensions follow formally.
\end{abstract}

\section*{Setup}

For integers $d \ge 1$ and $m \ge 2$, let
\[
\torus{d}{m}
=
\operatorname{Cay}\bigl((\mathbb Z/m\mathbb Z)^d,\{e_0,\dots,e_{d-1}\}\bigr).
\]
Thus $\torus{d}{m}$ has one outgoing edge in each coordinate direction. A
\emph{Hamilton decomposition} of $\torus{d}{m}$ is a partition of its arc set
into $d$ arc-disjoint directed Hamilton cycles.

\section*{The two-dimensional base case}

\begin{proposition}[Square lemma]
For every integer $n \ge 2$, the directed torus
\[
\torus{2}{n} \cong \cycle{n} \square \cycle{n}
\]
admits a Hamilton decomposition into two directed Hamilton cycles.
\end{proposition}

\begin{proof}
Work on the vertex set $\mathbb Z_n^2$. At each $(u,v)$ there are two outgoing
arcs:
\[
(u,v)\to(u+1,v),
\qquad
(u,v)\to(u,v+1).
\]
Define spanning subdigraphs $H_0$ and $H_1$ by
\[
H_0(u,v)=
\begin{cases}
(u+1,v), & u+v \not\equiv n-1 \pmod n,\\
(u,v+1), & u+v \equiv n-1 \pmod n,
\end{cases}
\]
and let $H_1$ be the complementary choice:
\[
H_1(u,v)=
\begin{cases}
(u,v+1), & u+v \not\equiv n-1 \pmod n,\\
(u+1,v), & u+v \equiv n-1 \pmod n.
\end{cases}
\]
These two rules clearly partition the two outgoing arcs at every vertex.

They also partition the two incoming arcs. Into $(u,v)$, the west incoming arc
comes from $(u-1,v)$ and is used by $H_0$ exactly when
\[
(u-1)+v \not\equiv n-1 \pmod n,
\]
that is, when $u+v \not\equiv 0 \pmod n$. The south incoming arc comes from
$(u,v-1)$ and is used by $H_0$ exactly when
\[
u+(v-1) \equiv n-1 \pmod n,
\]
that is, when $u+v \equiv 0 \pmod n$. So every vertex has indegree $1$ in
$H_0$, and therefore also in $H_1$. Hence both $H_0$ and $H_1$ are directed
$1$-factors.

Now conjugate $H_0$ by
\[
\Psi(u,v)=(s,t):=(u+v,v).
\]
Since $\Psi^{-1}(s,t)=(s-t,t)$, one checks directly that
\[
\Psi H_0 \Psi^{-1}(s,t)=
\begin{cases}
(s+1,t), & s \neq n-1,\\
(0,t+1), & s=n-1.
\end{cases}
\]
This is the standard odometer on $\mathbb Z_n^2$, so it is a single directed
cycle of length $n^2$. Hence $H_0$ is Hamilton.

Finally, $H_1$ is conjugate to $H_0$ by the coordinate swap
\[
(u,v)\longmapsto(v,u),
\]
so $H_1$ is Hamilton as well.
\end{proof}

\section*{Composite dimensions}

\begin{theorem}[Multiplicative closure]
Let $a,b \ge 2$. Assume that for every integer $n \ge 3$, both
\[
\torus{a}{n}
\qquad\text{and}\qquad
\torus{b}{n}
\]
admit Hamilton decompositions. Then for every integer $m \ge 3$, the directed
torus $\torus{ab}{m}$ also admits a Hamilton decomposition.
\end{theorem}

\begin{proof}
Fix $m \ge 3$, and write
\[
V:=((\mathbb Z/m\mathbb Z)^a)^b.
\]
By grouping the coordinates into $b$ consecutive blocks of size $a$, we may
identify
\[
\torus{ab}{m}
\cong
G^{(1)} \square G^{(2)} \square \cdots \square G^{(b)},
\]
where each factor $G^{(r)}$ is a copy of $\torus{a}{m}$.

For each factor $r \in \{1,\dots,b\}$, choose a Hamilton decomposition
\[
G^{(r)} = A_0^{(r)} \sqcup A_1^{(r)} \sqcup \cdots \sqcup A_{a-1}^{(r)},
\]
where each $A_i^{(r)}$ is a directed Hamilton cycle on the vertex set
$(\mathbb Z/m\mathbb Z)^a$. Every such cycle has length
\[
n := m^a.
\]

For each $i \in \{0,\dots,a-1\}$, define
\[
F_i := A_i^{(1)} \square A_i^{(2)} \square \cdots \square A_i^{(b)}.
\]
Since each $A_i^{(r)}$ is a directed cycle of length $n$, the product $F_i$ is
canonically isomorphic to
\[
\cycle{n}^{\square b} \cong \torus{b}{n}.
\]
By the hypothesis on dimension $b$, each $F_i$ therefore admits a Hamilton
decomposition into $b$ directed Hamilton cycles.

It remains to check that the subdigraphs $F_0,\dots,F_{a-1}$ partition the arc
set of $\torus{ab}{m}$. An arc of
\[
G^{(1)} \square \cdots \square G^{(b)}
\]
changes exactly one block, say the $r$-th block, and in that block it uses one
arc of $G^{(r)}$. Because the cycles
\[
A_0^{(r)},\dots,A_{a-1}^{(r)}
\]
partition the arc set of $G^{(r)}$, that changing-block arc belongs to a unique
$A_i^{(r)}$. By definition, the full product arc then belongs to the unique
subdigraph $F_i$.

Thus the $F_i$ are pairwise arc-disjoint and cover all arcs of $\torus{ab}{m}$.
Since each $F_i$ splits into $b$ Hamilton cycles, the whole digraph
$\torus{ab}{m}$ splits into
\[
a \cdot b = ab
\]
Hamilton cycles, as required.
\end{proof}

\begin{corollary}[Prime reduction]
Assume that for every prime $p$ and every integer $n \ge 3$, the directed torus
$\torus{p}{n}$ admits a Hamilton decomposition. Then for every integer $d \ge 2$
and every integer $m \ge 3$, the directed torus $\torus{d}{m}$ admits a
Hamilton decomposition.
\end{corollary}

\begin{proof}
Proceed by induction on $d$. If $d$ is prime, this is exactly the hypothesis.
If $d$ is composite, write $d=ab$ with $2 \le a,b < d$. By the induction
hypothesis, both $\torus{a}{n}$ and $\torus{b}{n}$ admit Hamilton
decompositions for every $n \ge 3$. The multiplicative-closure theorem then
gives a Hamilton decomposition of $\torus{d}{m}$ for every $m \ge 3$.
\end{proof}

\begin{corollary}[Immediate families from currently solved cases]
The following infinite families follow immediately.
\begin{enumerate}
\item Since $\torus{2}{n}$ is solved by the square lemma, any solved dimension
    $d$ automatically gives the doubled dimension $2d$.
\item Since $\torus{3}{n}$ is already solved in this repository for every
    $n \ge 3$, the dimensions $6=2\cdot 3$, $9=3\cdot 3$, $12=3\cdot 4$,
    $18=3\cdot 6$, and more generally every $2^{\alpha}3^{\beta}$ with
    $\alpha+\beta \ge 1$, are all forced by repeated product reduction.
\item Since $\torus{4}{n}$ is also solved here for every $n \ge 3$, the
    dimensions $8=2\cdot 4$, $16=4\cdot 4$, and every other dimension obtained
    by multiplying together already-solved dimensions are likewise forced.
\end{enumerate}
\end{corollary}

\begin{remark}
This argument is theorem-level rather than witness-level. It proves existence
for composite dimensions once smaller factors are known, but it does not
directly produce a tiny low-layer local rule table of the kind used in the
current explicit $d=3$ and $d=4$ constructions.
\end{remark}

\end{document}
